\documentclass[11pt,a4paper]{article}
\usepackage{fontspec}
\setmainfont{DejaVu Serif}[Scale=0.92]
\setsansfont{DejaVu Sans}[Scale=0.92]
\setmonofont{DejaVu Sans Mono}[Scale=0.85]
\usepackage{polyglossia}
\setmainlanguage{french}
\usepackage[margin=2.2cm]{geometry}
\usepackage{amsmath,amssymb}
\usepackage{booktabs,array,tabularx}
\usepackage[most]{tcolorbox}
\usepackage{xcolor}
\usepackage[hidelinks]{hyperref}
\usepackage{enumitem}
\setlist{itemsep=2pt,topsep=4pt}
\definecolor{bleu}{HTML}{1F5FA8}
\definecolor{vert}{HTML}{1D6B34}
\definecolor{rouge}{HTML}{9B2C2C}
\newtcolorbox{question}[1]{colback=blue!2,colframe=bleu!70,fonttitle=\bfseries,title={#1},breakable}
\newtcolorbox{reponse}{colback=gray!5,colframe=gray!40,boxrule=0.4pt,breakable}
\newfontfamily\symb{DejaVu Sans}
\newcommand{\ok}{\textcolor{vert}{{\symb ✔}}}
\newcommand{\ko}{\textcolor{rouge}{{\symb ✘}}}

\title{\textbf{Insurance Claims Assistance 19.5.0}\\[4pt]\large Démonstration commentée : cinq questions et le score des sinistres}
\date{Mesures effectuées le 18 novembre 2026 sur l'application en ligne}

\begin{document}
\sloppy
\maketitle

\section*{Ce que fait l'application}
L'application aide un gestionnaire de sinistres habitation. Elle n'automatise aucune décision.
\begin{enumerate}
\item \textbf{Assistant documentaire} : une question est rapprochée, par similarité sémantique
(embeddings bge-m3), des 201 documents d'assurance habitation. Les meilleurs documents sont lus
par un modèle de langage local (\texttt{llama3.2:3b} via Ollama), qui rédige une réponse courte
avec des extraits. \textbf{L'application vérifie elle-même} que chaque extrait figure mot pour mot
dans la source citée, masque toute phrase non appuyée, ajoute les conditions omises et décide du
statut affiché.
\item \textbf{Score des sinistres} : sur 120 dossiers \emph{fictifs}, une régression logistique
estime la probabilité d'acceptation d'un dossier et explique les facteurs qui la font monter ou
baisser ; des règles métier signalent les points bloquants et proposent une orientation.
\end{enumerate}

\section*{Comment lire une réponse}
\begin{tabularx}{\linewidth}{>{\bfseries}l X}
\toprule
Statut « réponse trouvée et justifiée » & le modèle conclut et ses extraits sont retrouvés dans les sources \\
Statut « réponse partielle » & la source ne couvre qu'une partie de la question, une phrase a été masquée ou une condition ajoutée \\
Statut « réponse absente des documents » & aucune source ne répond : rien n'est affirmé \\
Statut « non vérifiable » & la synthèse contenait une affirmation sans appui : seuls les extraits vérifiés sont montrés \\
{[S1], [S2]…} & renvoient aux sources listées sous la réponse \\
\ok\ / \ko & extrait retrouvé / non retrouvé dans la source indiquée \\
Similarité & proximité sémantique question–document (cosinus bge-m3), pas une probabilité d'exactitude \\
\bottomrule
\end{tabularx}

\section{Les cinq questions}
Chaque question a été posée une fois à l'application en ligne (réponses ci-dessous telles
qu'affichées). Les taux « mesuré » proviennent de 5 passages par question
(script \texttt{verifier\_application.py} du dossier \texttt{evaluation}).

%------------------------------------------------------------
\begin{question}{1. Réponse complète — « Que signifie louer un logement équipé de mobilier ? »}
Statut : \textbf{réponse trouvée et justifiée} ; 39 s ; mesuré 5/5.
\begin{reponse}
Louer un logement équipé de mobilier signifie que le propriétaire doit fournir un logement avec les
éléments indispensables pour y vivre normalement dès l'arrivée, tels que des sièges, des luminaires,
des étagères de rangement, etc. Le propriétaire doit également tenir un inventaire des équipements et
mobiliers mentionnant leur état d'usage. [S1]
\end{reponse}
Extraits : \ok\ « C'est un logement équipé d'éléments indispensables permettant d'y vivre
normalement dès l'arrivée dans les lieux. » [S1] ; \ko\ « Le propriétaire d'un meublé doit fournir
un logement habitable et paisible. » [S2].\\
Sources : S1 « Qu'est-ce qu'une location meublée ? », S2 « Assurance location meublée »…

\smallskip\textbf{Commentaire.} La question reformule le titre du document (« équipé de mobilier »
pour « meublée ») : c'est la recherche sémantique qui le retrouve en premier, sans mot commun.
Le modèle n'a lu que S1 (boucle « top-1 » : la meilleure source suffisait, réponse en 39 s au lieu
de 80 s). L'extrait attribué à S2 est marqué \ko\ : S2 n'a pas été transmise au modèle, la
citation est donc rejetée — la vérification fonctionne même quand la phrase est plausible.
\end{question}

%------------------------------------------------------------
\begin{question}{2. Réponse partielle avec restriction — « La moisissure sur mes murs est-elle prise en charge ? »}
Statut : \textbf{réponse partielle} ; 81 s ; mesuré 5/5.
\begin{reponse}
La moisissure sur les murs n'est pas prise en charge par l'assurance habitation MMA, car les
dommages causés par des champignons ou des moisissures ne sont pas assurés dans le cadre de la
garantie dégât des eaux. [S1] [S2]
\end{reponse}
Extraits : \ok\ « Dans le cadre de la garantie dégât des eaux, les dommages dus à l'humidité ou à la
condensation, les dommages causés par des champignons ou des moisissures, ne sont pas assurés »
[S2] ; \ok\ « Le contrat d'assurance habitation MMA propose le dépannage des installations de
serrurerie. » [S1].\\
Sources : S1 « … quelle prise en charge en cas de problème de serrure ? », S2 « Est-ce que les
dommages dus à l'humidité sont couverts ? »

\smallskip\textbf{Commentaire.} La conclusion est juste et appuyée mot pour mot par S2 : c'est
une \emph{exclusion} correctement trouvée. Le second extrait (serrurerie) est réel, donc
\ok, mais sans rapport : la vérification prouve qu'un extrait \emph{existe}, pas qu'il
\emph{justifie} la phrase. C'est la limite documentée qui a motivé l'essai du contrôle par inférence
textuelle (non retenu, voir l'addendum).
\end{question}

%------------------------------------------------------------
\begin{question}{3. Réponse partielle — « Mon emplacement de stationnement doit-il être assuré ? »}
Statut : \textbf{réponse partielle} ; 85 s ; mesuré 4/5.
\begin{reponse}
Si vous êtes propriétaire d'une copropriété, votre place de stationnement est souvent couverte par
l'assurance de la copropriété. [S1]\\
\emph{1 phrase de la synthèse du modèle ne figurait pas dans les sources citées et n'a pas été affichée.}
\end{reponse}
Extraits : \ok\ « Dans le cas d'une place de stationnement, tout particulièrement située au sein
d'une copropriété, vous n'avez pas besoin de souscrire une assurance parking. » ; \ok\ « La place de
stationnement est couverte par l'assurance de la copropriété, étant donné qu'elle est considérée
comme une partie commune… ».\\
Sources : S1 « L'assurance d'une place de stationnement », S2 « Faut-il assurer une place de parking ? »…

\smallskip\textbf{Commentaire.} Une phrase rédigée par le modèle n'était pas appuyée : elle est
\emph{masquée} et l'utilisateur en est informé. Le statut passe à « partielle » : la réponse est
vraie dans le cas de la copropriété, la source ne traite pas les autres cas. C'est le comportement
attendu : mieux vaut une réponse incomplète et sûre qu'une réponse complète inventée.
\end{question}

%------------------------------------------------------------
\begin{question}{4. Limite connue — « Quel est le délai de carence de mon assurance décès ? »}
Statut : \textbf{réponse partielle} (attendu : réponse absente) ; 82 s ; mesuré 0/5.
\begin{reponse}
Le délai de carence de votre assurance décès est de 90 jours après la prise d'effet de l'assurance
MMA de votre nouvelle habitation, sous condition que votre précédente habitation ait été également
assurée par MMA. [S1]
\end{reponse}
Extrait : \ok\ « les garanties souscrites pour la précédente habitation perdurent pendant 90 jours ».
Sources : S1 « Que couvre la garantie déménagement ? »…

\smallskip\textbf{Commentaire.} La question porte sur l'assurance décès, absente du corpus
(habitation uniquement). La recherche renvoie toujours les documents \emph{les plus proches},
même s'ils ne traitent pas le sujet ; le modèle a pris un vrai extrait sur le déménagement
(« 90 jours ») et l'a plaqué sur la question. Chaque mot vient de la question ou de la source :
le contrôle d'ancrage ne le voit pas. \textbf{C'est la limite mesurée} : sur 28 passages de
questions hors des documents, l'application ne s'abstient que 5 fois (18\,\%). Deux correcteurs
par inférence textuelle bloquaient 86\,\% de ces conclusions, mais aussi 7 bonnes réponses sur 15 :
ils n'ont pas été intégrés. La correction de fond — vérifier que le \emph{sujet} de la question
existe dans les documents avant de rédiger — est conçue dès le départ dans le projet suivant.
\end{question}

%------------------------------------------------------------
\begin{question}{5. « Je ne sais pas » étayé — bougie fondue sur la table de jardin}
Statut : \textbf{non vérifiable : extraits affichés} ; 59 s ; mesuré 2/5.
\begin{reponse}
La synthèse rédigée par le modèle contenait des affirmations qui ne figurent pas dans les sources
qu'il a citées. Pour ne rien affirmer sans preuve, l'application ne l'affiche pas : voici les
extraits vérifiés, à lire pour juger vous-même.
« Le contenu des boîtes aux lettres n'est pas assuré » [S1] ; « Les objets à risque de vol (voir
définition) appartenant à vos invités ne sont pas assurés » [S2].
\end{reponse}
Phrases bloquées : « Si vous avez acheté une assurance habitation, vous êtes couvert contre les
dégâts causés par les objets à risque de vol… » (termes absents des sources : \emph{acheté, dégâts,
causés}).

\smallskip\textbf{Commentaire.} Aucun document ne traite d'un objet fondu par la chaleur. Le
modèle a rédigé une conclusion sans rapport ; l'application l'a détectée (mots absents des
sources) et a basculé en mode secours. C'est le garde-fou central : \emph{aucune affirmation sans
preuve}. Sur 5 passages, ce garde-fou se déclenche 2 fois ; les 3 autres fois, le modèle conclut
avec des mots présents dans les sources, ce qui rejoint la limite de la question 4.
\end{question}

\section{Le score des sinistres (tableau de bord \texttt{/claims})}

\subsection{Indicateurs du portefeuille fictif}
120 dossiers : 72 acceptés, 27 refusés, 21 en attente. Le taux d'acceptation est calculé sur les
dossiers \emph{clos} seulement : $72/(72+27) = 72{,}73\,\%$ ; inclure les dossiers en attente le
ferait baisser artificiellement ($72/120 = 60\,\%$). Durée moyenne de traitement des dossiers clos :
22,1 jours (médiane 23).

\subsection{Le modèle}
Régression logistique sur 7 variables (type de sinistre, garantie, canal de déclaration, dossier
complet, montant demandé, ancienneté du client, sinistres précédents) ; variables catégorielles en
indicatrices, variables numériques standardisées. Pour le dossier $i$ :
\[
\operatorname{logit}(p_i)=\ln\frac{p_i}{1-p_i}=\beta_0+\sum_{f}\beta_f\,x_{i,f},
\qquad p_i=\frac{1}{1+e^{-(\beta_0+\sum_f\beta_f x_{i,f})}} ,
\]
où $p_i$ est la probabilité d'acceptation, $x_{i,f}$ la valeur (standardisée ou indicatrice) de la
variable $f$ et $\beta_f$ son coefficient. La \emph{contribution} affichée d'un facteur est
$\beta_f\,x_{i,f}$ : positive, elle fait monter le score ; négative, elle le fait baisser.

\begin{center}
\begin{tabularx}{\linewidth}{lr X}
\toprule
Mesure (25 dossiers de test) & Valeur & Signification \\
\midrule
Exactitude & 68\,\% (17/25) & IC de Wilson 95\,\% [48\,\% ; 83\,\%] \\
Exactitude équilibrée & 60,3\,\% & moyenne des rappels des deux classes (corrige le déséquilibre 73/27) \\
F1 (classe « acceptée ») & 77,8\,\% & $2PR/(P+R)$, précision $P$ et rappel $R$ \\
ROC AUC & 75,4\,\% & probabilité qu'un dossier accepté ait un score supérieur à un refusé \\
\bottomrule
\end{tabularx}
\end{center}

\textbf{À dire :} avec 25 dossiers de test, l'intervalle de l'exactitude va de 48\,\% à 83\,\% ;
sur des données fictives, ces chiffres prouvent que la chaîne fonctionne, pas que le modèle est
bon. L'exactitude équilibrée (60\,\%) montre que les refus sont mal reconnus : un modèle qui
accepterait tout aurait déjà 73\,\% d'exactitude. D'où la règle : \emph{aucune décision
automatique}.

\subsection{Un dossier expliqué : C00020 (catastrophe naturelle, en attente)}
Probabilité estimée : 21,3\,\%, soit $\operatorname{logit}=\ln(0{,}213/0{,}787)=-1{,}31$.

\begin{center}
\begin{tabular}{lr}
\toprule
Facteur & Contribution $\beta_f x_{i,f}$ \\
\midrule
Type de sinistre = Catastrophe naturelle & $+0{,}771$ \\
Dossier complet = Oui & $+0{,}643$ \\
Ancienneté du client & $+0{,}195$ \\
Garantie = À vérifier & $-1{,}235$ \\
Canal de déclaration = Courtier & $-0{,}363$ \\
Sinistres précédents & $-0{,}189$ \\
\bottomrule
\end{tabular}
\end{center}
La somme des six facteurs affichés vaut $-0{,}18$ ; la constante $\beta_0$ et les autres variables
apportent le reste, jusqu'à $-1{,}31$. Le facteur dominant est « garantie à vérifier ».

\textbf{Règles métier} (indépendantes du score) : point \emph{bloquant} « garantie à vérifier » ;
alerte « en attente depuis 318 jours : anomalie probable, à escalader ». \textbf{Orientation} :
« compléter le dossier avant décision », car un point bloquant rend le score inexploitable seul.
Les clauses du contrat correspondant au type de sinistre sont proposées, chacune avec un lien vers
l'assistant documentaire : les deux briques de l'application se rejoignent ici.

Pour comparaison, C00001 (bris de glace, accepté) : probabilité 85,4\,\%, garantie « Oui », dossier
complet ; validation humaine toujours requise.

\section{Questions probables et réponses courtes}
\begin{description}
\item[Pourquoi un modèle local ?] Les données de sinistres ne quittent pas l'infrastructure ;
coût fixe ; même modèle en ligne et en local.
\item[Comment savez-vous que la nouvelle recherche est meilleure ?] Bon document parmi les 3
premiers : 26/30 contre 13/30 sur des questions tenues à part (McNemar, $p=0{,}001$).
\item[Et de bout en bout ?] 41/90 réponses conformes contre 28/90 ; 8 questions améliorées, 2 en
recul ; $p=0{,}11$ : gain net, pas encore significatif avec 30 questions.
\item[Que se passe-t-il si le serveur tombe ?] Même image Docker en local (trois commandes), avec
adresse publique temporaire par Cloudflare Tunnel ; procédure vérifiée sur une machine vierge.
\item[Quelle est la principale limite ?] Les questions hors du périmètre (question 4) ; elle est
mesurée, et la correction testée a été écartée parce qu'elle faisait perdre de bonnes réponses.
\end{description}

\end{document}
